% Laget automatisk av elevark/lag_elevark.py fra opplegg/faraday-lenz.html. Endringer her blir overskrevet.
\documentclass[11pt]{article}
\usepackage{elevark}
\begin{document}
\elevarkhode{Fysikk 2 · Elektromagnetisk induksjon}{Faradays og Lenz' lov – arbeidsark}{https://fysikklab.no/sim/magnet-og-spole.html}{Magnet og spole}

\section*{Oppgaver}

Oppgave 1–5 bruker Magnet og spole\footnote{\url{fysikklab.no/sim/magnet-og-spole.html}}. Oppgave 6–8 er uten simulering.

\begin{oppgaveliste}
\oppgave \fase{gjett}{Gjett}Magneten ligger i ro midt inne i spolen. Blir det indusert en spenning? Dra magneten dit og sjekk.
\svar
\oppgave \fase{test}{Test}Trykk Start. Når er spenningen størst, og når er den null? Sammenlign med fluksen (den stiplede kurven).
\svar
\oppgave \fase{test}{Undersøk}Hvordan endrer toppverdien av spenningen seg når du dobler farten? Når du dobler antall vindinger?
\svar
\oppgave \fase{forklar}{Forklar}Se på polene i spolen når magneten nærmer seg og når den fjerner seg. Forklar med Lenz' lov hvorfor spolen alltid motvirker bevegelsen.
\svar
\oppgave \fase{forklar}{Forklar}Snu magneten. Hva skjer med fortegnet på spenningen og med polene? Hvorfor?
\svar
\oppgave \fase{test}{Regn}En kvadratisk sløyfe med sider 10~cm står i et felt på 0,20~T. Normalen danner 60° med feltet. Regn ut fluksen. Sløyfa dreies til normalen er parallell med feltet på 0,10~s. Hvor stor er den gjennomsnittlige induserte spenningen?
\svar
\oppgave \fase{test}{Regn}En metallstang på 0,50~m glir med 3,0~m/s på to skinner, vinkelrett på et felt på 0,40~T. Skinnene er koblet sammen med en motstand på 2,0~Ω. Finn den induserte spenningen, strømmen og kraften som trengs for å holde farten konstant.
\svar
\oppgave \fase{utfordring}{Utfordring}Vis at effekten du må tilføre stanga i oppgave 7 $(F \cdot v)$ er lik den elektriske effekten i motstanden. Hva forteller det om Lenz' lov?
\svar
\end{oppgaveliste}

\end{document}
